\section{Labeled graph models}

\subsection{A canonical set-based Petersen model}

A coordinate-free labeled realization of the Petersen graph is the Kneser
graph

\[
P=KG(5,2).
\]

Its vertices are the two-element subsets of

\[
[5]=\{1,2,3,4,5\},
\]

and two vertices are adjacent when they are disjoint.

Thus

\[
V(P)
=
\{
12,13,14,15,23,24,25,34,35,45
\},
\]

where \(ij\) abbreviates \(\{i,j\}\).

Its fifteen edges are

\[
\begin{aligned}
&
\{12,34\},\{12,35\},\{12,45\},\\
&
\{13,24\},\{13,25\},\{13,45\},\\
&
\{14,23\},\{14,25\},\{14,35\},\\
&
\{15,23\},\{15,24\},\{15,34\},\\
&
\{23,45\},\{24,35\},\{25,34\}.
\end{aligned}
\]

These fifteen edges form the vertex set of

\[
Y=L(P).
\]

Two vertices of \(Y\) are adjacent when the corresponding Petersen edges
share an endpoint.

\subsection{The double-cover model}

The vertices of \(X\) are

\[
(e,\varepsilon),
\qquad
e\in E(P),
\quad
\varepsilon\in\Ftwo.
\]

Adjacency is

\[
(e,\varepsilon)
\sim
(f,\varepsilon+1)
\]

precisely when \(e\) and \(f\) share a Petersen vertex.

This set-based model is sufficient for all conceptual proofs in the paper.

\subsection{Certificate integer labels}

The machine-readable certificates use integer vertex labels. Let
\(e_b\) denote the Petersen edge in row \(b\) of
Table~\ref{tab:certificate-label-map}. The certificate convention is

\[
(e_b,\varepsilon)
\longmapsto
2b+\varepsilon,
\qquad
\varepsilon\in\Ftwo.
\]

\begin{table}[h]
\centering
\begin{tabular}{rrrr}
\toprule
Base index \(b\) & Petersen edge \(e_b\) &
\((e_b,0)\) & \((e_b,1)\)\\
\midrule
0 & $\{12,34\}$ & 0 & 1\\
1 & $\{12,45\}$ & 2 & 3\\
2 & $\{23,45\}$ & 4 & 5\\
3 & $\{15,23\}$ & 6 & 7\\
4 & $\{15,34\}$ & 8 & 9\\
5 & $\{14,23\}$ & 10 & 11\\
6 & $\{14,25\}$ & 12 & 13\\
7 & $\{13,25\}$ & 14 & 15\\
8 & $\{13,45\}$ & 16 & 17\\
9 & $\{25,34\}$ & 18 & 19\\
10 & $\{14,35\}$ & 20 & 21\\
11 & $\{12,35\}$ & 22 & 23\\
12 & $\{15,24\}$ & 24 & 25\\
13 & $\{24,35\}$ & 26 & 27\\
14 & $\{13,24\}$ & 28 & 29\\
\bottomrule
\end{tabular}
\caption{Exact conversion between the set-based cover coordinates and the
integer labels used by the machine-readable certificates.}
\label{tab:certificate-label-map}
\end{table}

The underlying Project 41 alignment is recorded in

\begin{verbatim}
sources/project41-paper42/a5_v4_k22_four_slot_alignment_audit_019.json
\end{verbatim}

with SHA-256 digest

\begin{verbatim}
e8bbfe48053bdc1a80407ce7a26bf02b12f83bfcf695fed0073a9a8581ae896a
\end{verbatim}

Audit \(019\) fixes the fifteen base indices and their Petersen-edge
alignment. It uses an integer-labeled standard Petersen model. The
Project 42 bridge certificate

\begin{verbatim}
artifacts/json/project42_kneser_label_bridge_045.json
\end{verbatim}

enumerates all \(120\) graph isomorphisms from that standard model to the
fixed \(KG(5,2)\) model used in this appendix and selects the
lexicographically least image tuple. Its SHA-256 digest is

\begin{verbatim}
381d0b5e0ef0a11fa152417f3e6b4a74b0e38a449f12d06e9a1616ea2e2943c6
\end{verbatim}

The choice of this bridge is editorial coordinate data. It changes no
abstract graph, covering, automorphism-group, or quotient-frame claim.
